\documentclass{beamer}

\pdfmapfile{+sansmathaccent.map}


\mode<presentation>
{
	\usetheme{Warsaw} % or try Darmstadt, Madrid, Warsaw, Rochester, CambridgeUS, ...
	\usecolortheme{seahorse} % or try seahorse, beaver, crane, wolverine, ...
	\usefonttheme{serif}  % or try serif, structurebold, ...
	\setbeamertemplate{navigation symbols}{}
	\setbeamertemplate{caption}[numbered]
} 


%%%%%%%%%%%%%%%%%%%%%%%%%%%%
% itemize settings


%%%%%%%%%%%%%%%%%%%%%%%%%%%%
% itemize settings

\definecolor{myhotpink}{RGB}{255, 80, 200}
\definecolor{mywarmpink}{RGB}{255, 60, 160}
\definecolor{mylightpink}{RGB}{255, 80, 200}
\definecolor{mypink}{RGB}{255, 30, 80}
\definecolor{mydarkpink}{RGB}{155, 25, 60}

\definecolor{mypaleblue}{RGB}{240, 240, 255}
\definecolor{mylightblue}{RGB}{120, 150, 255}
\definecolor{myblue}{RGB}{90, 90, 255}
\definecolor{mygblue}{RGB}{70, 110, 240}
\definecolor{mycourseblue}{RGB}{120, 140, 255}
\definecolor{mycourselightblue}{RGB}{245, 250, 255}
\definecolor{mydarkblue}{RGB}{0, 0, 180}
\definecolor{myblackblue}{RGB}{40, 40, 120}

\definecolor{mygreen}{RGB}{0, 200, 0}
\definecolor{mygreen2}{RGB}{245, 255, 230}
\definecolor{mydarkgreen}{RGB}{0, 120, 0}
\definecolor{mydarkgreen}{RGB}{0, 120, 0}


\definecolor{mygray}{gray}{0.8}
\definecolor{mydarkgray}{RGB}{80, 80, 160}
\definecolor{mylightgray}{RGB}{160, 160, 160}

\definecolor{mydarkred}{RGB}{160, 30, 30}
\definecolor{mylightred}{RGB}{255, 150, 150}
\definecolor{myred}{RGB}{200, 110, 110}
\definecolor{myblackred}{RGB}{120, 40, 40}

\definecolor{mypink}{RGB}{255, 30, 80}
\definecolor{myhotpink}{RGB}{255, 80, 200}
\definecolor{mywarmpink}{RGB}{255, 60, 160}
\definecolor{mylightpink}{RGB}{255, 80, 200}
\definecolor{mydarkpink}{RGB}{155, 25, 60}
\definecolor{mywhitepink}{RGB}{255, 240, 240}

\definecolor{mydarkcolor}{RGB}{60, 25, 155}
\definecolor{mylightcolor}{RGB}{130, 180, 250}



\setbeamertemplate{itemize items}[default]

\setbeamertemplate{itemize item}{\color{myblackblue}$\blacksquare$}
\setbeamertemplate{itemize subitem}{\color{mygblue}$\blacktriangleright$}
\setbeamertemplate{itemize subsubitem}{\color{mygray}$\blacksquare$}

\setbeamercolor{palette quaternary}{fg=white,bg=mycourseblue}
\setbeamercolor{titlelike}{parent=palette quaternary}

\setbeamercolor{palette quaternary2}{fg=black,bg=mycourselightblue}
\setbeamercolor{frametitle}{parent=palette quaternary2}

\setbeamerfont{frametitle}{size=\Large,series=\scshape}
\setbeamerfont{framesubtitle}{size=\normalsize,series=\upshape}





%%%%%%%%%%%%%%%%%%%%%%%%%%%%
% block settings

\setbeamercolor{block title}{bg=red!30,fg=black}

\setbeamercolor*{block title example}{bg=mygreen!40!white,fg=black}

\setbeamercolor*{block body example}{fg= black, bg= mygreen2}


%%%%%%%%%%%%%%%%%%%%%%%%%%%%
% URL settings
\hypersetup{
	colorlinks=true,
	linkcolor=blue,
	filecolor=blue,      
	urlcolor=blue,
}

%%%%%%%%%%%%%%%%%%%%%%%%%%

\renewcommand{\familydefault}{\rmdefault}


\newcommand{\mycolor}[1] {
\setbeamercolor{normal text}{fg={#1}} % Recolors standard body text
\setbeamercolor{structure}{fg={#1}}   % Recolors bullets and structure items
\usebeamercolor[fg]{normal text}      % Applies the normal text color immediately
}


\usepackage{amsmath}
\usepackage{mathtools}

\usepackage{cancel}
\usepackage{centernot}

\usepackage{changepage}
\usepackage{subcaption}

\usepackage{qrcode}

\newcommand{\bo}[1] {\mathbf{#1}}
\newcommand{\R}{\mathbb{R}} 
\newcommand{\T}{^\top}     



\newcommand{\mydate}{Fall 2026}

\newcommand{\mygit}{\textcolor{blue}{\href{https://github.com/SergeiSa/Advanced-Control-2026}{github.com/SergeiSa/Advanced-Control-2026}}}

\newcommand{\myqr}{ \textcolor{black}{\qrcode[height=1.5in]{https://github.com/SergeiSa/Advanced-Control-2026}}
}

\newcommand{\myqrframe}{
	\begin{frame}
		\centerline{Lecture slides are available via Github:}
		\bigskip
		\centerline{\mygit}
		\bigskip
		\myqr
	\end{frame}
}


\newcommand{\bref}[2] {\textcolor{blue}{\href{#1}{#2}}}

%%%%%%%%%%%%%%%%%%%%%%%%%%%%
% code settings

\usepackage{listings}
\usepackage{color}
% \definecolor{mygreen}{rgb}{0,0.6,0}
% \definecolor{mygray}{rgb}{0.5,0.5,0.5}
\definecolor{mymauve}{rgb}{0.58,0,0.82}
\lstset{ 
	backgroundcolor=\color{white},   % choose the background color; you must add \usepackage{color} or \usepackage{xcolor}; should come as last argument
	basicstyle=\footnotesize,        % the size of the fonts that are used for the code
	breakatwhitespace=false,         % sets if automatic breaks should only happen at whitespace
	breaklines=true,                 % sets automatic line breaking
	captionpos=b,                    % sets the caption-position to bottom
	commentstyle=\color{mygreen},    % comment style
	deletekeywords={...},            % if you want to delete keywords from the given language
	escapeinside={\%*}{*)},          % if you want to add LaTeX within your code
	extendedchars=true,              % lets you use non-ASCII characters; for 8-bits encodings only, does not work with UTF-8
	firstnumber=0000,                % start line enumeration with line 0000
	frame=single,	                   % adds a frame around the code
	keepspaces=true,                 % keeps spaces in text, useful for keeping indentation of code (possibly needs columns=flexible)
	keywordstyle=\color{blue},       % keyword style
	language=Octave,                 % the language of the code
	morekeywords={*,...},            % if you want to add more keywords to the set
	numbers=left,                    % where to put the line-numbers; possible values are (none, left, right)
	numbersep=5pt,                   % how far the line-numbers are from the code
	numberstyle=\tiny\color{mygray}, % the style that is used for the line-numbers
	rulecolor=\color{black},         % if not set, the frame-color may be changed on line-breaks within not-black text (e.g. comments (green here))
	showspaces=false,                % show spaces everywhere adding particular underscores; it overrides 'showstringspaces'
	showstringspaces=false,          % underline spaces within strings only
	showtabs=false,                  % show tabs within strings adding particular underscores
	stepnumber=2,                    % the step between two line-numbers. If it's 1, each line will be numbered
	stringstyle=\color{mymauve},     % string literal style
	tabsize=2,	                   % sets default tabsize to 2 spaces
	title=\lstname                   % show the filename of files included with \lstinputlisting; also try caption instead of title
}


%%%%%%%%%%%%%%%%%%%%%%%%%%%%
% URL settings
\hypersetup{
	colorlinks=false,
	linkcolor=blue,
	filecolor=blue,      
	urlcolor=blue,
}

%%%%%%%%%%%%%%%%%%%%%%%%%%

%%%%%%%%%%%%%%%%%%%%%%%%%%%%
% tikz settings

\usepackage{tikz}
\tikzset{every picture/.style={line width=0.75pt}}
\usetikzlibrary{shapes.geometric, decorations.markings, positioning, patterns, calc, arrows.meta}


\subtitle{Advanced Control}
\author{by Sergei Savin}
\centering
\date{\mydate}



\title{Bellman Equation, LQR}
\begin{document}
\maketitle


\begin{frame}{Content}
	\begin{itemize}
		\item Optimality and the Value Function
		\item Bellman Equation
		\item Finite-Horizon Discrete-Time LQR
		\begin{itemize}
			\item Optimal controller
			\item Riccati recursion
		\end{itemize}
		\item Infinite-Horizon Discrete-Time LQR
		\begin{itemize}
			\item Algebraic Riccati equation
			\item Solving the Riccati equation numerically
		\end{itemize}
	\end{itemize}
\end{frame}



%------------------------------------------------
\begin{frame}{Optimality}
	\begin{flushleft}
		
		Consider a discrete-time dynamical system:
		\[
		x_{k+1}=f(x_k,u_k).
		\]
		
		A control law is given by
		\[
		u_k=\pi(x_k).
		\]
		
		For a finite time horizon of $N$ steps, let the cost be
		\[
		J=\phi(x_N)+\sum_{k=0}^{N-1}\ell(x_k,u_k).
		\]
		
		\medskip
		
		We can now pose the following question:
		
		\medskip
		
		\emph{What is the optimal control law $\pi^*$ and the
			corresponding trajectory that minimize the cost $J$?}
		
	\end{flushleft}
\end{frame}
%------------------------------------------------

\begin{frame}{Value Function}
	\begin{flushleft}
		
		We define the value function $V_k(x)$ as the minimum
		cost achievable when starting from $x_k=x$ and
		considering the remaining $N-k$ steps:
		\[
		\begin{aligned}
			V_k(x)=\min_{u_k,\ldots,u_{N-1}}\quad
			&\phi(x_N)+\sum_{i=k}^{N-1}\ell(x_i,u_i)\\
			\text{subject to}\quad
			&x_{i+1}=f(x_i,u_i),\\
			&x_k=x.
		\end{aligned}
		\]
		
		\medskip
		
		The value function represents the \textbf{optimal
			cost-to-go}: the lowest possible cost for a fixed
		number of remaining time steps.
		
	\end{flushleft}
\end{frame}
%------------------------------------------------

\begin{frame}{Bellman Equation}
	\begin{flushleft}
		
		The Bellman equation describes how the value function
		at one step is related to the value function at the
		next step:
		\[
		\boxed{
			V_k(x)=\min_u
			\left[
			\ell(x,u)+V_{k+1}\bigl(f(x,u)\bigr)
			\right].
		}
		\]
		
		\medskip
		
		The optimal cost-to-go consists of two parts:
		\begin{itemize}
			\item The immediate cost $\ell(x,u)$.
			\item The optimal cost-to-go from the next state.
		\end{itemize}
		
		\medskip
		
		At the terminal step, no time remains, so the value
		function equals the terminal cost:
		\[
		V_N(x)=\phi(x).
		\]
		
	\end{flushleft}
\end{frame}
%------------------------------------------------




%------------------------------------------------
\begin{frame}{Linear Case}
	\begin{flushleft}
		
		Consider linear dynamics:
		\[
		x_{k+1}=Ax_k+Bu_k.
		\]
		
		For now, suppose that the control law is linear:
		\[
		u_k=-Kx_k.
		\]
		
		The closed-loop dynamics are
		\[
		x_{k+1}=(A-BK)x_k.
		\]
		
		\medskip
		
		We consider a quadratic cost function:
		\[
		J=x_N^\top Q_Nx_N+
		\sum_{i=0}^{N-1}
		\left(x_i^\top Qx_i+u_i^\top Ru_i\right),
		\]
		where $Q_N\succeq0$, $Q\succeq0$, and $R\succ0$.
		
	\end{flushleft}
\end{frame}
%------------------------------------------------

\begin{frame}{Linear Case: Value Function}
	\begin{flushleft}
		{\small
		Let us first consider a fixed linear controller
		$u_k=-Kx_k$.
		
		Under this controller, the future state depends on the
		current state as follows:
		\[
		x_i=(A-BK)^{i-k}x_k.
		\]
		
		Moreover,
		\[
		x_i^\top Qx_i+u_i^\top Ru_i
		=x_i^\top(Q+K^\top RK)x_i.
		\]
		
		\medskip
		
		Therefore, the cost-to-go under this controller is
		quadratic:
		\[
		V_k(x)=x^\top S_kx,
		\]
		where
		\[
		\begin{aligned}
			S_k={}&
			\bigl((A-BK)^{N-k}\bigr)^\top
			Q_N(A-BK)^{N-k}\\
			&+\sum_{i=k}^{N-1}
			\bigl((A-BK)^{i-k}\bigr)^\top
			(Q+K^\top RK)
			(A-BK)^{i-k}.
		\end{aligned}
		\]
		
		\medskip
		
		The same quadratic form will describe the optimal
		value function, but with matrices $S_k$ determined
		by the Bellman equation.}
		
	\end{flushleft}
\end{frame}
%------------------------------------------------



%------------------------------------------------
\begin{frame}{Linear Case: Bellman Equation}
	\begin{flushleft}
		
		The Bellman equation is
		\[
		V_k(x)=\min_u
		\left[
		\ell(x,u)+V_{k+1}(Ax+Bu)
		\right].
		\]
		
		For the quadratic cost, suppose that
		\[
		V_k(x)=x^\top S_kx.
		\]
		
		Then the Bellman equation becomes
		\[
		x_k^\top S_kx_k
		=\min_{u_k}\left[
		x_k^\top Qx_k+u_k^\top Ru_k
		+x_{k+1}^\top S_{k+1}x_{k+1}
		\right].
		\]
		
		\medskip
		
		Since $x_{k+1}=Ax_k+Bu_k$, we can write, using
		$x=x_k$ and $u=u_k$:
		\[
		\boxed{
			x^\top S_kx=\min_u
			\left[
			x^\top Qx+u^\top Ru
			+(Ax+Bu)^\top S_{k+1}(Ax+Bu)
			\right].
		}
		\]
		
	\end{flushleft}
\end{frame}
%------------------------------------------------

\begin{frame}{Optimal Controller}
	\begin{flushleft}
		
		To find the optimal control, we differentiate the
		expression inside the minimization with respect to $u$
		and set the derivative to zero:
		\[
		\begin{aligned}
			&2Ru+2B^\top S_{k+1}Bu +2B^\top S_{k+1}Ax=0.
		\end{aligned}
		\]
		
		Equivalently,
		\[
		(R+B^\top S_{k+1}B)u
		=-B^\top S_{k+1}Ax.
		\]
		
		\medskip
		
		Solving for $u$ gives the optimal control law:
		\[
		\boxed{
			u=-K_kx,
		}
		\]
		where
		\[
		\boxed{
			K_k=(R+B^\top S_{k+1}B)^{-1}
			B^\top S_{k+1}A.
		}
		\]
		
		\medskip
		
		To compute the optimal gain at each step, we need
		the matrix $S_{k+1}$.
		
	\end{flushleft}
\end{frame}
%------------------------------------------------

\begin{frame}{Discrete-Time Riccati Recursion}
	\begin{flushleft}
		
		Substituting the optimal control into the Bellman
		equation gives
		\[
		\begin{aligned}
			x^\top S_kx
			={}&x^\top Qx+x^\top A^\top S_{k+1}Ax\\
			&-x^\top A^\top S_{k+1}B (R+B^\top S_{k+1}B)^{-1} B^\top S_{k+1}Ax.
		\end{aligned}
		\]
		
		\medskip
		
		Since this equality must hold for every $x$, the
		matrices of the quadratic forms must be equal.
		We obtain the discrete-time Riccati recursion:
		\[
		\boxed{
			\begin{aligned}
				S_k={}&Q+A^\top S_{k+1}A-A^\top S_{k+1}B
				(R+B^\top S_{k+1}B)^{-1} B^\top S_{k+1}A
			\end{aligned}
		}
		\]
		
		\medskip
		
		The terminal condition follows from
		$V_N(x)=x^\top Q_Nx$:
		\[
		S_N=Q_N.
		\]
		
		Starting from $S_N$, we compute $S_{N-1}$,
		then $S_{N-2}$, and continue backward until $S_0$.
		
	\end{flushleft}
\end{frame}
%------------------------------------------------




%------------------------------------------------
\begin{frame}{Infinite-Horizon Cost}
	\begin{flushleft}
		
		Consider a different cost function:
		\[
		J=\sum_{k=0}^{\infty}
		\left(x_k^\top Qx_k+u_k^\top Ru_k\right).
		\]
		
		\medskip
		
		Notice two important differences from the finite-horizon
		case:
		\begin{itemize}
			\item There is no terminal cost.
			\item The cost is accumulated over an infinite
			number of time steps.
		\end{itemize}
		
		\medskip
		
		This leads to the \textbf{infinite-horizon
			discrete-time LQR} problem.
		
	\end{flushleft}
\end{frame}
%------------------------------------------------

\begin{frame}{Infinite-Horizon Case: Value Function}
	\begin{flushleft}
		
		In the finite-horizon case, the value function $V_k(x)$
		depends on the number of time steps remaining.
		
		\medskip
		
		In the infinite-horizon case, there are always infinitely
		many steps remaining. For a time-invariant system and
		cost, the optimal cost-to-go does not depend on the
		current time step.
		
		\medskip
		
		We can therefore write
		\[
		V(x)=\min_{u_0,u_1,\ldots}
		\sum_{k=0}^{\infty}
		\left(x_k^\top Qx_k+u_k^\top Ru_k\right).
		\]
		
		\medskip
		
		As before, the value function is quadratic:
		\[
		\boxed{V(x)=x^\top Sx.}
		\]
		
		Here, $S$ is a constant matrix rather than a
		time-dependent sequence of matrices $S_k$.
		
	\end{flushleft}
\end{frame}
%------------------------------------------------

\begin{frame}{Infinite-Horizon Discrete-Time LQR}
	\begin{flushleft}
		
		The Bellman equation becomes
		\[
		V(x)=\min_u
		\left[
		x^\top Qx+u^\top Ru+V(Ax+Bu)
		\right].
		\]
		
		Substituting $V(x)=x^\top Sx$, we obtain
		\[
		x^\top Sx=\min_u
		\left[
		x^\top Qx+u^\top Ru
		+(Ax+Bu)^\top S(Ax+Bu)
		\right].
		\]
		
		\medskip
		
		Repeating the steps from the finite-horizon case gives
		the optimal control law:
		\[
		\boxed{
			u=-Kx,
			\qquad
			K=(R+B^\top SB)^{-1}B^\top SA.
		}
		\]
		
		\medskip
		
		Substituting the optimal control into the Bellman
		equation yields the discrete-time algebraic
		Riccati equation:
		\[
		\boxed{
			\begin{aligned}
				S={}&Q+A^\top SA\\
				&-A^\top SB(R+B^\top SB)^{-1}B^\top SA.
			\end{aligned}
		}
		\]
		
	\end{flushleft}
\end{frame}
%------------------------------------------------

\begin{frame}{Solving the Riccati Equation}
	\begin{flushleft}
		
		The discrete-time algebraic Riccati equation can be
		solved numerically using standard control libraries.
		
		\medskip
		
		\begin{itemize}
			\item \textbf{Python (SciPy):}
			\texttt{scipy.linalg.solve\_discrete\_are(A, B, Q, R)}
			
			\medskip
			
			\item \textbf{MATLAB:}
			\texttt{dlqr(A, B, Q, R)}
			
			\medskip
			
			\item \textbf{Julia:}
			\texttt{lqr(Discrete, A, B, Q, R)}
		\end{itemize}
		
		\medskip
		
		Once $S$ is obtained, the optimal feedback gain is
		\[
		K=(R+B^\top SB)^{-1}B^\top SA.
		\]
		
		The resulting control law $u=-Kx$ is time-invariant.
		
	\end{flushleft}
\end{frame}
%------------------------------------------------



\begin{frame}{Literature}
	\begin{itemize}
		
		\item MIT OpenCourseWare:
		\bref{https://ocw.mit.edu/courses/16-323-principles-of-optimal-control-spring-2008/pages/lecture-notes/}
		{Principles of Optimal Control: Lecture 3, Dynamic Programming and Discrete LQR.}
		
		\item Russ Tedrake -- MIT:
		\bref{https://underactuated.mit.edu/lqr.html}
		{Underactuated Robotics: Linear Quadratic Regulators, including the discrete-time Riccati equation.}
		
		\item University of Washington:
		\bref{https://faculty.washington.edu/ratliffl/teaching/2020_LinearSystems.pdf}
		{Lecture 10: Linear Quadratic Regulator -- Discrete Time.}
		
		\item SciPy documentation:
		\bref{https://docs.scipy.org/doc/scipy/reference/generated/scipy.linalg.solve_discrete_are.html}
		{Linear Algebra: \texttt{solve\_discrete\_are}, numerical solution of the discrete-time algebraic Riccati equation.}
		
	\end{itemize}
\end{frame}





\myqrframe




\end{document}
